\section{The family and its complete cycle census}\label{sec:main}

An explicit H\'enon map can have a large rational periodic set while
leaving its exact cycle structure undetermined. For the discrete-sine
family, the essential distinction is between the small-amplitude core
and a narrow boundary on which new periods occur. We determine both
pieces, with their actual point labels and ordinary iteration times.

For an odd integer $d=2k+1\ge3$, define
\begin{equation}\label{eq:family}
 s_d(Y)=\sum_{j=0}^{k}(-1)^{k-j}
 \frac{Y\prod_{i=1}^{j}(Y^2-i^2)}{(2j+1)!},
 \qquad h_d(x,y)=(y,-x+s_d(y)).
\end{equation}
The empty product is $1$. The degree of $s_d$ is exactly $d$, and
$h_d^{-1}(x,y)=(s_d(x)-y,x)$; thus $h_d$ is a polynomial automorphism
of $\Q^2$ with Jacobian determinant $1$. Throughout, a primitive cycle
means an orbit of least positive period, counted once under cyclic
rotation. We do not identify a cycle with its negative or its reversal.

Put
\begin{equation}\label{eq:radius}
 R=\frac{d+1}{2}=3q+s,\qquad s\in\{0,1,2\}.
\end{equation}
Here $q$ is an integer quotient, not a finite-field cardinality. The
allowed ranges are $q\ge1$ when $s=0,1$, and $q\ge0$ when $s=2$.
Let $C_n(d)$ count primitive rational cycles of length $n$.

\begin{theorem}[Complete rational cycle classification]\label{thm:main}
For the map \eqref{eq:family}, every rational periodic point is integral
and lies in $[-R-2,R+2]^2$. All primitive rational cycles are given by
the following three additive tables. An omitted period has multiplicity
zero; entries of the same period in different tables are added.

\begin{center}
\begin{minipage}{\linewidth}
\centering
\begin{tabular}{ll}
\toprule
Condition & Central cycles, written as period: multiplicity\\
\midrule
$R$ odd & $1:1,\quad 5:2,\quad 6:1$\\
$R$ even & $1:1,\quad 3:2,\quad 10:1$\\
\bottomrule
\end{tabular}
\captionof{table}{The central part, containing $17$ points in either phase.}
\label{tab:central}
\end{minipage}
\end{center}

\begin{center}
\begin{minipage}{\linewidth}
\centering
\begin{tabular}{cccc}
\toprule
$s$ & Bulk $C_4$ & Bulk $C_{12}$ & Bulk $C_{20}$\\
\midrule
$0$ & $q(q+1)$ & $q(q-1)$ & $q(q-1)$\\
$1$ & $q(q+1)$ & $q(q+1)$ & $q(q-1)$\\
$2$ & $q(q+1)$ & $q(q+1)$ & $q(q+1)$\\
\bottomrule
\end{tabular}
\captionof{table}{Generic core cycles, excluding the central part.}
\label{tab:bulk-main}
\end{minipage}
\end{center}

\begin{center}
\begin{minipage}{\linewidth}
\centering
\begin{tabular}{cl}
\toprule
$s$ & Boundary cycles, written as period: multiplicity\\
\midrule
$0$ & $4:2q+2,\quad 14:2$\\
$1$ & $6:2,\quad 16:2,\quad 36:q-1,\quad (16q+6):1$\\
$2$ & $4:2,\quad 6:1,\quad 8:2q$\\
\bottomrule
\end{tabular}
\captionof{table}{All cycles meeting the boundary section.}
\label{tab:boundary-main}
\end{minipage}
\end{center}

The tables specify an explicit point classification, not only
multiplicities. The positive-phase core consists exactly of the clipped
residue rectangles in Table~\ref{tab:affine}, including its central
exceptions. The negative-phase point labels are obtained by the
conjugacy of Proposition~\ref{prop:phase}. The boundary cycles are the
complete first-return families of Section~\ref{sec:closure}; their
intermediate points are recovered by iterating \eqref{eq:family}.
\end{theorem}

The word ``bulk'' below always refers to the generic core cycles in
Table~\ref{tab:bulk-main}; it does not include the $17$ central points.
The word ``boundary'' refers to an orbit that meets the section, not to
the assertion that all its points lie on the boundary of a square.

\begin{corollary}[Point totals and ordinary return law]\label{cor:counts}
The total number $P(d)=\#\Per(h_d,\Q)$ is
\begin{equation}\label{eq:total}
P(d)=
\begin{cases}
36q^2-20q+53=(3d^2-4d+152)/3,& R=3q,\quad d\equiv5\pmod6,\\
36q^2+48q+31=d^2+6d+24,& R=3q+1,\quad d\equiv1\pmod6,\\
36q^2+52q+31=d^2+8d/3+14,& R=3q+2,\quad d\equiv3\pmod6.
\end{cases}
\end{equation}
For each $n\ge1$, let $N_n(d)=\#\Fix(h_d^n,\Q)$. Then
\begin{equation}\label{eq:zeta}
 N_n(d)=\sum_{\ell\mid n}\ell C_\ell(d),\qquad
 Z_d(u)=\exp\left(\sum_{n\ge1}\frac{N_n(d)}n u^n\right)
       =\prod_{\ell\ge1}(1-u^\ell)^{-C_\ell(d)}.
\end{equation}
The product is finite, and the exponential identity is an identity of
formal power series, as well as an analytic identity for $|u|<1$.
\end{corollary}

The proof proceeds by integrality and escape, exact clipping of all
affine core cells, and a signed first-return calculation on two boundary
layers. Keeping the return signs is essential: a sign-negative return
cycle of total time $T$ produces ordinary period $2T$, not $T$.
